\documentclass[11pt,a4paper]{article}
\usepackage{fontspec}
\setmainfont{Arial}
\setsansfont{Arial}
\setmonofont{Arial}
\usepackage{polyglossia}
\setdefaultlanguage{russian}
\setotherlanguage{english}
\usepackage{geometry,amsmath,amssymb,booktabs,longtable,tabularx,array,enumitem,xcolor,hyperref,fancyhdr,tikz}
\usetikzlibrary{arrows.meta,positioning,shapes.geometric}
\geometry{margin=18mm}
\hypersetup{colorlinks=true,linkcolor=blue!55!black,urlcolor=blue!65!black}
\setlist{nosep,leftmargin=*}
\setlength{\parindent}{0pt}
\setlength{\parskip}{5pt}
\definecolor{navy}{HTML}{17365D}\definecolor{teal}{HTML}{167D8D}\definecolor{pale}{HTML}{EEF5F7}\definecolor{warn}{HTML}{8A4B08}
\setlength{\headheight}{14pt}
\pagestyle{fancy}\fancyhf{}\lhead{MATEMAFIJA}\rhead{Team Project Guide}\cfoot{\thepage}
\newcommand{\term}[1]{\textbf{#1}}
\newcommand{\code}[1]{\texttt{#1}}
\newcommand{\boxtext}[1]{\vspace{2mm}\noindent\fcolorbox{teal}{pale}{\parbox{0.94\linewidth}{#1}}\vspace{2mm}}
\title{\vspace{-1cm}\textbf{MATEMAFIJA}\\[3mm]\Large Полный team guide по проекту и защите\\[2mm]\large Hack4Vilnius EV Charging Policy Lab}
\author{Рабочий документ Глеба и Галины}
\date{Финальная hackathon-версия 3.0 --- 10 октября 2026}

\begin{document}
\maketitle
\boxtext{\textbf{Главная мысль проекта.} MATEMAFIJA --- это не обещание, что новая цена или credit обязательно устранит дефицит. Это \textbf{policy stress-testing digital twin}: инструмент, который на реальных OCPI/DATEX данных показывает, при каких явно заданных price-and-credit response assumptions balanced contribution-to-credit policy проходит system, redistribution, user-proxy и accounting gates, и какой пилот нужен городу для проверки предположений.}
\tableofcontents
\newpage

\section{Как пользоваться guide}
Сначала прочитайте разделы 2--5: после них вы сможете за 10 минут объяснить идею. Затем изучите 6--13 для понимания реализации. Перед защитой используйте разделы 14--17: demo, pitch, вопросы жюри и шпаргалка.

\textbf{Разделение ролей при подготовке:} Глеб ведёт data pipeline, model contracts, demo и технические ответы; Галина ведёт problem story, users/business, city benefit и основную речь. Но каждый должен уметь восстановить весь flow по памяти, объяснить provenance каждого числа и назвать claim boundary.

\section{Проект за 10 минут}
\subsection{Проблема}
Зарядная инфраструктура может периодически иметь нулевую доступность. Но в имеющихся данных нет физической очереди, времени ожидания, момента окончания зарядки или времени, когда водитель остаётся после окончания. Новая политика ещё не внедрялась, поэтому её реальный поведенческий эффект тоже неизвестен.

\subsection{Что мы сделали}
Мы организовали реальные OCPI и DATEX наблюдения, исправили временную семантику, восстановили интервалы состояний EVSE, построили time series и метрики scarcity, воспроизвели фактический OCPI тариф F1 и определили OSA-v1 scarcity signal. Поверх него спроектирован balanced contribution-to-credit mechanism: prior funded MCU поддерживают выбор qualified receiver sites. Fixed и adaptive сценарии получают один baseline и одинаковые случайные числа; результат --- joint observable robustness map, а не единичный оптимистичный forecast.

\subsection{Что получает город}
Город получает ответ не «какую цену или награду немедленно поставить», а: \emph{при каких price-and-credit response assumptions политика проходит system, redistribution, user-proxy и accounting gates; и как спроектировать ограниченный измеримый пилот}.

\subsection{Кто пользователь и кто платит}
Ежедневный beneficiary --- водитель, которому показывается пригодная альтернатива и bounded incentive. Operational user --- charging operator, который тестирует network utilisation, receiver capacity и fundability. Первый реалистичный customer --- партнёрство города и CPO для policy analysis и bounded pilot. На MVP нельзя обещать готовую revenue model: pilot должен определить funding source, operator liability, acceptance и settlement. Product сегодня --- decision-support и pilot-design tool, а не consumer wallet.

\section{Архитектура целиком}
\begin{center}
\begin{tikzpicture}[node distance=5mm, every node/.style={font=\small}, b/.style={draw,rounded corners,fill=pale,text width=12.8cm,align=center,minimum height=8mm}, arr/.style={-{Stealth},thick,teal}]
\node[b] (raw) {OCPI + DATEX: locations, EVSE, connectors, statuses, tariffs, timestamps};
\node[b,below=of raw] (db) {SQLite v2: source-native records, snapshots, relationships, provenance};
\node[b,below=of db] (state) {Reconstruction: state intervals $\rightarrow$ flow/historical series $\rightarrow$ scarcity and derivative statistics};
\node[b,below=of state] (base) {Baseline: F1 observed tariff replay + empirical transition hazards + observed/bootstrap path};
\node[b,below=of base] (policy) {OSA-v1 signal + contribution pool + funded MCU alternatives};
\node[b,below=of policy] (resp) {Assumed response: release, accepted-start, credit/distance choice, delay};
\node[b,below=of resp] (sim) {Paired Monte Carlo + safety + append-only MCU/capacity ledger};
\node[b,below=of sim] (out) {System + redistribution + user-proxy + accounting $\rightarrow$ joint observable robustness};
\foreach \a/\b in {raw/db,db/state,state/base,base/policy,policy/resp,resp/sim,sim/out} \draw[arr] (\a)--(\b);
\end{tikzpicture}
\end{center}

\textbf{Логика сообщений между блоками:} raw data сообщает факты; reconstruction сообщает допустимые признаки; policy сообщает scarcity signal и MCU rules; response model сообщает условный механизм поведения; simulator сообщает возможное состояние и ledger events; evaluation сообщает раздельные domain metrics. Ни один блок не должен переименовывать допущение в наблюдение.

\section{Данные и происхождение переменных}
\begin{tabularx}{\textwidth}{>{\bfseries}p{2.6cm}X p{4.0cm}}
\toprule Класс & Значение & Пример \\\midrule
observed & непосредственно получено от источника & OCPI status, timestamp, tariff element \\
reconstructed & детерминированно восстановлено из наблюдений & interval, availability ratio, sustained duration \\
estimated & статистически оценено при выполнении quality gate & empirical transition hazard \\
policy-defined & нормативный параметр решения & grace=15 min, cap=0.05 EUR/min \\
assumed & неизвестная гипотеза сценария & price/credit sensitivities, delay, receiver choice \\
simulated & результат модели, не реальный факт & counterfactual EVSE counts \\
\bottomrule
\end{tabularx}

\textbf{OCPI} даёт структуру locations--EVSE--connectors, текущие статусы и компонентные тарифы. \textbf{DATEX II} даёт независимую статическую/динамическую информацию и corroboration. Источники сохраняются раздельно; cross-source mapping обязан иметь confidence и review status.

Все timestamps хранятся как timezone-aware UTC; календарные признаки рассчитываются в \code{Europe/Vilnius}. Наивный OCPI timestamp интерпретируется как локальное время Вильнюса и затем переводится в UTC. Collection gap не считается нулевым спросом.

\boxtext{\textbf{Текущий снимок БД (10.10.2026, около 17:47 EEST):} 2,836,399 OCPI EVSE-status observations; 1,984 current OCPI location records; 1,707 tariffs; 5,900 tariff assignments; 4.08 дня временного окна. В записях с city, содержащим Vilnius: 892 locations и 2,120 EVSE, из них 1,331 available и 180 charging в текущем snapshot. Это source records, не утверждение об уникальных физических станциях. Formal 7-day gate ещё не выполнен, поэтому formal effectiveness heatmap заблокирован. Для demo разрешён отдельно маркированный illustrative response preview: он строится только из declared assumptions, не участвует в robustness decision и не является evidence или forecast. Числа пересчитать перед записью: collector продолжает работу.}

\section{Термины без жаргона}
\begin{longtable}{p{3.5cm}p{12.5cm}}
\toprule Термин & Простое значение \\\midrule\endhead
station/site/location & физическая точка или площадка зарядки; OCPI Location --- основной site ID MVP.\\
EVSE & управляемая зарядная единица; именно EVSE является единицей capacity, не connector.\\
connector & физический разъём EVSE; число разъёмов не равно числу одновременных сессий.\\
snapshot & снимок состояний в конкретный момент.\\
state interval & реконструированный отрезок времени, на котором состояние не менялось.\\
proxy & измеримая замена ненаблюдаемому понятию; не тождество.\\
scarcity & нехватка доступной operational capacity; не наблюдаемая очередь.\\
baseline & фактическая/эмпирическая точка сравнения до новой policy.\\
hazard & условная интенсивность перехода состояния за следующий шаг времени.\\
counterfactual & «что было бы при другой policy»; здесь это сценарий, не факт.\\
$\lambda$ & сила adaptive price adjustment от 0 до 1.\\
$\kappa$ & доля EVSE/sites, охваченная политикой.\\
$R_i(t)$ & rollout indicator: применяется ли policy на site $i$ в момент $t$.\\
$\theta_r$ & assumed sensitivity release transition к цене.\\
$\theta_a$ & assumed sensitivity accepted-start transition к цене.\\
$\tau$ & задержка реакции.\\
Monte Carlo & повторение стохастического сценария для оценки распределения результата.\\
common random numbers & одинаковые случайные числа в paired arms; уменьшают шум сравнения.\\
Wilson interval & устойчивый доверительный интервал для доли успешных replications.\\
robustness region & область assumptions, где нижняя граница вероятности улучшения проходит порог.\\
spatial spillover & scarcity, перенесённая на соседние stations из-за displaced use.\\
MCU & integer mobility-credit unit для scenario accounting; не money и не DAEI.\\
observable policy robust & joint system/redistribution/user-proxy/accounting success в пределах доступных evidence.\\
\bottomrule
\end{longtable}

\section{От snapshot к scarcity}
Для site $i$ и времени $t$:
\[
A_i(t)=\frac{\text{available EVSE}}{\text{effective operational capacity}},\qquad
C_i(t)=1-A_i(t).
\]
$C_i$ --- operational scarcity ratio. Нулевая доступность допустима как evidence только при положительной capacity, достаточном coverage и пригодном snapshot. Equipment unavailability считается отдельно: сломанная инфраструктура не должна автоматически запускать ценовой механизм.

Ключевые outputs: mean availability, operational scarcity, share пригодных snapshots с zero availability, длительность эпизодов (только для stateful/observed trajectories), release/start proxy rates. Мы никогда не называем $A=0$ очередью.

\section{F1: реальный fixed baseline}
\term{F1 observed tariff replay} воспроизводит структурированный OCPI tariff: \code{ENERGY}, \code{TIME}, \code{PARKING\_TIME}, \code{FLAT}, ограничения и VAT/currency. DATEX pricing используется как corroboration, а не как скрытая замена OCPI.

Нельзя свести любой тариф к одной цене EUR/min. Полная стоимость требует Session/CDR/meter quantities: kWh, charging duration, parking duration. Если их нет, механизм тарифа показывается, но monetary total честно withheld.

\section{OSA-v1: адаптивная policy}
OSA-v1 сохраняет идею pressure + grace + nonlinear adjustment, но заменяет ненаблюдаемые queue $Q$ и post-charge idle time на доступные признаки.
\[
P_i(t)=\operatorname{clip}\!\left(\frac{C_i(t)-C_{trigger}}{1-C_{trigger}},0,1\right),
\quad
R_i(t)=\operatorname{clip}\!\left(\frac{D_i(t)-T_{grace}}{T_{ramp}},0,1\right)
\]
\[
S_i(t)=P_i(t)^{\gamma}R_i(t),\qquad
\Delta p_{raw}(i,t)=\Delta p_{max}S_i(t).
\]
Здесь $D_i$ --- reconstructed lower-bound duration устойчивого scarcity condition, не idle time автомобиля. Она строится по упорядоченной последовательности snapshot конкретной станции, а не приходит готовым полем из OCPI. Для каждого snapshot проверяется одновременно:
\begin{enumerate}
\item snapshot пригоден для capacity evidence: coverage не ниже 0.80, capacity положительна и availability определена;
\item $C_i(t)>C_{trigger}$, то есть operational pressure действительно выше порога;
\item хотя бы один EVSE имеет статус \code{CHARGING}: существует наблюдаемый proxy для возможного release channel;
\item доля equipment-unavailable EVSE не выше 0.50, чтобы поломка не интерпретировалась как demand scarcity;
\item разрыв от предыдущего snapshot не превышает две ожидаемые cadence (в MVP: $2\times5=10$ min).
\end{enumerate}

Время первого snapshot непрерывной qualifying-последовательности обозначим $t_i^{start}$. Тогда
\[
D_i(t)=\frac{t-t_i^{start}}{60\ \text{seconds}}
\]
в минутах. На первом qualifying snapshot $D_i=0$. Если следующий snapshot также проходит все условия, $D_i$ растёт на фактически прошедшее время. Любой failed condition, ineligible snapshot или gap больше 10 минут разрывает эпизод: $t_i^{start}$ забывается, а $D_i$ становится 0 до начала новой qualifying-последовательности. Поэтому величина является \emph{нижней границей} продолжительности: дефицит мог начаться между двумя наблюдениями, но мы считаем его только с первого подтверждённого snapshot.

Пример при cadence 5 min: snapshots в 18:00, 18:05, 18:10 и 18:15 проходят условия, поэтому $D_i=0,5,10,15$. В 18:20 $D_i=20$, и persistence factor становится $(20-15)/30=1/6$. Если запись 18:25 потеряна, но 18:30 получена, gap равен 10 min и последовательность ещё сохраняется; если следующий snapshot только в 18:31, gap 11 min превышает guard, и новый эпизод начинается с $D_i=0$. Это защищает policy от ложной уверенности при дырках в данных.

Параметры MVP: $C_{trigger}=0.50$, $T_{grace}=15$ min, $T_{ramp}=30$ min, $\gamma=1.2$, $\Delta p_{max}=0.05$ EUR/min; coverage gate 0.80; equipment guard 0.50. Это versioned test policy, не оптимальные значения для Вильнюса.

Итоговый driver реакции:
\[
d_i(t)=R_i^{rollout}(t)\,\lambda\,\Delta p_{raw}(i,t).
\]
Не путайте persistence $R$ в формуле OSA с rollout $R_i^{rollout}$: в речи называйте их \emph{persistence factor} и \emph{rollout indicator}.

Policy возвращает \code{active}, \code{inactive} или \code{withheld}. Withheld не равен нулю: это отсутствие достаточного основания для решения.

\section{Три режима simulator}
\begin{itemize}
\item \code{observed\_replay}: применяет pricing mechanism к реальному пути без поведенческой реакции; разрешённый claim --- \emph{mechanism only}.
\item \code{assumption\_driven\_counterfactual}: меняет transitions по явно заданным response parameters; разрешённый claim --- \emph{scenario conditional}.
\item \code{validated\_counterfactual}: будущий режим после внешней или пилотной валидации; разрешён ограниченный \emph{counterfactual effect}.
\end{itemize}

\section{Baseline transition model и trajectories}
Из qualified history оцениваются два proxy-перехода: \code{charging $\to$ available} (release) и \code{available $\to$ charging} (accepted start). Они stratified по weekday/hour с fallback. Это не arrivals, queue, charge completion или post-charge idle.

Два генератора baseline: observed-path replay и local-day block bootstrap. Bootstrap сохраняет суточную структуру, station heterogeneity, autocorrelation и совместность состояний лучше, чем независимая генерация рядов.

Readiness gates: descriptive --- минимум 2 eligible snapshots; historical --- минимум 7 дней, 7 local dates и 80\% expected coverage; transition --- historical gate, минимум 40 complete intervals и 20 transitions каждого направления; F1 --- все current connectors имеют разрешимый tariff. Counterfactual требует transition + F1.

\section{Поведенческие response channels}
Release channel:
\[
h^{release}_{1}=h^{release}_{0}\exp\{\theta_r d_i(t-\tau)\}.
\]
При $\theta_r>0$ price signal повышает вероятность proxy release.

Arrival/accepted-start channel:
\[
h^{start}_{1}=h^{start}_{0}\exp\{-\theta_a d_i(t-\tau)\}.
\]
Он уменьшает accepted starts. Поэтому улучшение scarcity нельзя показывать без displaced/deferred start proxy: иначе модель может «победить», просто подавив использование.

$\theta_r,\theta_a,\tau$ имеют provenance \code{assumed}. При нулевом driver adaptive arm обязан совпасть с fixed arm (identity test).

\section{Simulation, MCU circulation, evaluation и robustness}
\subsection{Что именно моделируется}
Раздел объединяет две связанные, но логически раздельные части. \term{Upstream simulation} отвечает на вопрос, как assumed response на policy signal может изменить proxy-переходы \code{charging $\to$ available} и \code{available $\to$ charging}. \term{Downstream credit circulation} отвечает на вопрос, что произойдёт с displaced/deferred start candidates после появления funded alternatives: будет ли выбран receiver, можно ли безопасно зарезервировать capacity и MCU, и сохраняется ли accounting identity.

Полная причинно-вычислительная цепочка сценария:
\[
\begin{aligned}
&\boxed{\text{observed/reconstructed baseline}}
\to\boxed{\text{F1 + OSA signal}}
\to\boxed{\text{assumed release/start response}}\\
&\hspace{8mm}\to\boxed{\text{redistribution candidates}}
\to\boxed{\text{funded receiver offers}}\\
&\hspace{8mm}\to\boxed{\text{choice + safety + ledger}}
\to\boxed{\text{four-domain evaluation}}.
\end{aligned}
\]
Каждый блок сохраняет provenance. Первые данные observed/reconstructed; policy parameters policy-defined; behavioural coefficients assumed; trajectories, choices и ledger events simulated. Благодаря этому итог нельзя ошибочно прочитать как prediction реального поведения.

\subsection{Paired experiment и one-step projection}
Upstream grid варьирует $\lambda\times\theta_r\times\theta_a\times\tau$; downstream grid добавляет credit/distance sensitivities, reserve и residual receiver capacity. Fixed/adaptive получают один trajectory и common random numbers. Funded offers проходят safety, choice, atomic MCU/capacity reservation и append-only ledger replay. One-step origin projection:
\[
A'=A+\text{releases}-\text{starts},\qquad C'=C-\text{releases}+\text{starts}.
\]
Это условная one-step модель, не полный forecast demand. Поскольку projections пока не fed forward, episode-duration metrics для них отключены.

\term{Paired} означает, что fixed и adaptive arms получают один исходный baseline и одинаковые random uniforms. Если outcome различается, причина --- policy/response parameters, а не независимый случайный шум. При $\lambda=0$ или нулевой response sensitivity adaptive arm обязан дать identity с fixed arm; это проверяется тестами.

Переменные $A'$ и $C'$ не становятся новым наблюдаемым состоянием станции. Они показывают только ближайшее условное следствие дополнительных release и изменённых accepted starts. Модель не переносит этот projection автоматически в следующий timestamp; поэтому она не заявляет полноценный demand forecast и не вычисляет counterfactual episode duration.

Для каждой метрики считается $\Delta M=M_{adaptive}-M_{fixed}$. Для zero-availability и scarcity отрицательное значение хорошо; для availability --- положительное.

Каждый run оценивается отдельно по system, redistribution, user-proxy и accounting domains. Joint run success требует прохождения всех четырёх gates; missing evidence остаётся \code{not\_evaluable}. Scenario получает \code{observable\_policy\_robust} только при минимум 20 valid replications, доле non-evaluable не более 0.20 и нижней 95\% Wilson bound joint success не меньше 0.80. Это не validated welfare/fairness и не pilot readiness.

Четыре домена нельзя взаимозаменять:
\begin{itemize}
\item \textbf{System}: изменение zero-availability share, mean availability, operational scarcity, release и displaced-start proxies относительно F1.
\item \textbf{Redistribution}: сколько candidates получили funded offer, выбрали receiver или остались outside; сколько capacity потреблено и сколько offers заблокировал residual-capacity guard.
\item \textbf{User-proxy}: размер предложенных/выбранных MCU и straight-line distance. Total session cost, road travel time, welfare и demographic fairness остаются \code{not\_evaluable}.
\item \textbf{Accounting}: funded, available, reserved, outstanding и settled MCU, issuance/redemption и нулевая reconciliation difference.
\end{itemize}
Joint success является конъюнкцией: хороший system effect не может компенсировать перенос дефицита на receiver, плохой user proxy или сломанный ledger. Wilson lower bound добавляет защиту от вывода по случайно удачной малой выборке: важна не только observed success share, но и её консервативная нижняя граница.

\subsection{Contribution и MCU accounting}
MCU --- целое scenario-accounting unit: $1$ MCU имеет policy-defined reference $0.01$ EUR-equivalent eligible charging service. Это не деньги, balance пользователя, cashback или DAEI conversion. В режиме \code{unit\_accounting} реальная session duration не выдумывается; используется reference interval $q_{ref}=5$ min:
\[
c_i^{MCU}(t)=\operatorname{round\_half\_up}\!\left(d_i(t)\,q_{ref}\,100\right).
\]
При cap $0.05$ EUR/min максимум --- 25 MCU на qualifying origin opportunity. Quote не является funding. Только completed simulated origin service может создать future inflow; current offer финансируется prior pool.

Ledger invariant после каждого события:
\[
funded=available+reserved+outstanding+settled.
\]
Reserve и offerable amount:
\[
reserve_{target}=\left\lceil r(available+reserved+outstanding)\right\rceil,
\qquad offerable=\max(0,available-reserve_{target}).
\]
Offer не двигает MCU; reservation переводит available $\to$ reserved; issuance --- reserved $\to$ outstanding; redemption --- outstanding $\to$ settled. MCU и receiver capacity резервируются атомарно.

Это различие устраняет double spending. Например, при 100 available MCU и reserve target 10 MCU максимально offerable 90 MCU. Показ пользователю offer на 20 MCU ещё не уменьшает balance. После успешного выбора ledger одновременно резервирует 20 MCU и одну capacity unit receiver: состояния становятся 80 available и 20 reserved. Только simulated receiver start переводит 20 reserved в 20 outstanding и потребляет capacity. Позднее redemption переводит outstanding в settled. Если service не начинается, reservation освобождается, а 20 MCU и capacity возвращаются; частичного движения только одного ресурса быть не может.

\subsection{Credit-sensitive receiver choice}
Для opportunity $o$ и qualified receiver $j$:
\[
V_{oj}=\beta_c c_{oj}-\beta_d d_{oj}+\beta_a A_j,
\qquad
P(j\mid o)=\frac{e^{V_{oj}}}{1+\sum_{k\in J_o}e^{V_{ok}}}.
\]
Outside option имеет utility 0 и вероятность $1/(1+\sum e^{V})$. Здесь $c$ --- funded MCU offer, $d$ --- straight-line distance, $A$ --- receiver availability; $\beta$ имеют provenance \code{assumed}. Distance не равен route time, а probability не является observed choice share.

Знаки коэффициентов задают прозрачную гипотезу: больший funded credit повышает utility, большая distance снижает её, а availability влияет напрямую только если интерфейс действительно показывает её пользователю. По умолчанию direct availability effect выключен, чтобы availability не действовала скрыто дважды --- как eligibility guard и как perceived benefit. Вектор вероятностей нормализуется вместе с обязательным outside outcome. Поэтому модель никогда не принуждает candidate выбрать станцию: при непривлекательных или отсутствующих offers probability outside остаётся положительной или равной 1.

Stochastic draw означает только simulated selection. Затем safety layer заново проверяет expiry, MCU cap, offerable pool, distance, текущую receiver availability, наличие capacity и minimum residual capacity. Selection становится reservation только после успешной атомарной проверки. Иными словами:
\[
\begin{aligned}
\text{offer shown}&\ne\text{receiver selected}\ne\text{capacity reserved}\\
&\ne\text{service started}\ne\text{MCU issued}\ne\text{MCU redeemed}.
\end{aligned}
\]
Эти состояния намеренно разделены и записываются append-only events; replay должен после каждого события восстановить тот же balance и lifecycle.

\section{Spatial redistribution}
\subsection{Зачем нужен пространственный слой}
Displaced/deferred accepted starts не могут просто исчезнуть из системы. Иначе scarcity можно было бы искусственно «улучшить», подавив starts на origin station и не показав, куда делась потребность. Пространственный слой сохраняет каждый integer candidate и допускает только два terminal outcomes: выбор qualified receiver или \code{outside/unassigned}. Последнее не называется отказом, потерянной поездкой или неудовлетворённым реальным demand: текущие данные этого не наблюдают.

Candidate появляется upstream только из парного различия accepted-start proxy между fixed и adaptive arms. Это не человек и не индивидуальная сессия, а единица simulated displaced-use accounting. Origin, timestamp, scenario и replication сохраняются, поэтому downstream результат можно точно связать с породившим его upstream run.

\subsection{Receiver eligibility и capacity}
Потенциальный receiver должен отличаться от origin, иметь координаты, находиться внутри policy-defined distance gate и обладать доступной operational capacity. Haversine distance по координатам --- straight-line proxy; она не утверждает road distance или travel time. Без данных нельзя гарантировать connector/power compatibility, accessibility, cross-operator acceptance или полную стоимость сессии, поэтому эти ограничения остаются явными warnings.

Usable receiving capacity ограничивается консервативно. В legacy spatial allocator:
\[
K_j^{usable}(t)=\left\lfloor A^{count}_j(t)\,f_{capacity}\right\rfloor,
\]
где $A^{count}_j$ --- projected available EVSE count, а $f_{capacity}\in[0,1]$ --- assumed доля, разрешённая для redistribution. В credit-circulation path safety guard дополнительно требует, чтобы после reservation оставалась минимум одна residual capacity unit. Поэтому улучшение origin не должно немедленно создавать zero availability на receiver.

\subsection{Два реализованных механизма}
Репозиторий сохраняет два совместимых уровня:
\begin{enumerate}
\item \textbf{Legacy capacity-constrained allocator}: применяет seeded participation probability, строит origin--receiver edges, сортирует их глобально nearest-first и распределяет integer starts без overbooking. Он полезен как простой network stress test.
\item \textbf{Current credit-sensitive path}: получает immutable funded offers, вычисляет choice probabilities с explicit outside option, применяет safety guards и только затем атомарно резервирует receiver capacity и MCU. Именно этот путь соответствует новой contribution-to-credit концепции.
\end{enumerate}
Nearest-first не используется как доказательство пользовательского выбора в credit-sensitive path: там distance является лишь одним assumed utility component. Оба механизма seeded и воспроизводимы, но отвечают на разные вопросы --- физически возможное распределение против гипотетического выбора funded alternative.

\subsection{Conservation и пример}
Для каждого origin и каждого decision timestamp должны выполняться identities:
\[
N_{candidate}=N_{receiver\ selected}+N_{outside},
\qquad
N_{displaced}=N_{allocated}+N_{unassigned}.
\]
Для receiver:
\[
N_{allocated,j}\le K_j^{usable},
\qquad
A^{after}_j=A^{before}_j-N_{allocated,j}\ge K_j^{residual}.
\]
Например, три candidates рассматривают receiver с двумя available units. Если policy требует оставить одну residual unit, безопасно зарезервировать можно только одну. Даже если второй candidate выбирает тот же receiver, повторная pre-reservation проверка видит обновлённый ledger/capacity state и блокирует overbooking; outcome остаётся outside/unassigned либо выбирается другая безопасная alternative. Итоги обязаны показать одну realised redistribution, два outside outcomes и одну consumed capacity unit, а не скрыть заблокированные попытки.

\subsection{Что разрешено утверждать}
Обязательно показываются receiver/outside outcomes, capacity consumed, residual-capacity blocks, distance proxies и ledger reconciliation. Разрешённый вывод: \emph{при данных assumptions сценарий перераспределяет такую долю proxy candidates, не нарушая заданные capacity и accounting guards}. Нельзя утверждать, что реальные водители поедут по этим маршрутам, что travel time приемлем, что individual welfare вырос или что достигнута demographic fairness. Для таких выводов нужны pilot choice records, Session/CDR, road routing, compatibility и privacy-safe user/access evidence.

\section{Карта репозитория}
\begin{longtable}{p{5.3cm}p{10.7cm}}
\toprule Путь & Роль \\\midrule\endhead
\code{src/collect.py} & цикл сбора; \code{--once}, dashboard mode.\\
\code{src/ocpi.py}, \code{datex.py} & source clients/parsing.\\
\code{src/database\_v2.py}, \code{schema\_v2.sql} & SQLite schema, relationships, writes/reads.\\
\code{src/historical.py}, \code{flow.py} & state intervals и time series.\\
\code{src/scarcity.py}, \code{derivative\_statistics.py} & evidence-based metrics.\\
\code{src/baseline\_readiness.py} & gates и причины ineligibility.\\
\code{src/transition\_model.py}, \code{trajectory.py} & hazards и baseline generators.\\
\code{src/fixed\_policy.py} & F1 tariff replay.\\
\code{src/adaptive\_policy.py} & OSA-v1 mechanism.\\
\code{src/policy\_integration.py} & F1 + adjustment + rollout/intensity.\\
\code{src/simulation\_contract.py} & типы, provenance, invariants.\\
\code{src/release\_response.py}, \code{arrival\_response.py} & behavioural assumptions.\\
\code{src/simulation\_grid.py} & paired Monte Carlo runner.\\
\code{src/simulation\_evaluation.py}, \code{robustness.py} & metrics, Wilson region, thresholds.\\
\code{src/receiver\_choice.py}, \code{simulation\_ledger.py} & credit-sensitive choice и append-only MCU lifecycle.\\
\code{src/fairness\_safety.py} & pre-reservation guards и distributional proxy audit.\\
\code{src/credit\_simulation\_grid.py} & downstream paired credit/choice/ledger grid.\\
\code{src/credit\_simulation\_evaluation.py}, \code{credit\_robustness.py} & four-domain evaluation и joint observable robustness.\\
\code{src/spatial\_redistribution.py} & network spillover prototype.\\
\code{src/dashboard.py} & raw evidence dashboard.\\
\code{src/policy\_dashboard.py}, \code{policy\_lab.py} & Policy Lab UI.\\
\code{docs/} & формальные contracts и claim boundaries.\\
\code{tests/} & 156 automated tests на текущем состоянии.\\
\bottomrule
\end{longtable}

Команды:
\begin{verbatim}
python src/collect.py --once
python src/collect.py
python -m streamlit run src/dashboard.py
python -m streamlit run src/policy_dashboard.py
python -m unittest discover -s tests -v
\end{verbatim}

\section{Что уже есть и чего ещё нет}
\textbf{Есть:} нормализованная БД, timezone rules, intervals/time series, scarcity и derivative statistics, readiness gates, F1, OSA-v1, typed response, receiver choice, MCU ledger, safety guards, paired upstream/downstream grids, four-domain evaluation, joint observable robustness и six-tab Policy Lab.

\textbf{Ещё нет как доказанного результата:} фактической elasticity водителей; observed queue/waiting/post-charge idle; causal proof; Session/CDR totals; road travel time; validated welfare/fairness; operator settlement; production deployment. Receiver eligibility и global offer allocation остаются approved design contracts и explicit inputs. Эти пункты --- честная граница MVP и программа пилота.

\section{Какие outputs можно показать}
\begin{longtable}{p{4.0cm}p{7.2cm}p{4.5cm}}
\toprule Output & Что показываем & Разрешённый смысл \\\midrule\endhead
Live evidence & карта, sources, timestamps, status counts, tariff components & реальные current observations \\
Readiness & history span, coverage и blocker reasons & почему model output разрешён или withheld \\
OSA-v1 & pressure, persistence, bounded adjustment & policy mechanism, не фактический bill \\
Response curves & release/start probability и receiver choice & sensitivity assumption, не measured elasticity \\
Illustrative preview & assumed response strength $\times$ policy intensity & mechanism only; не evidence/forecast \\
MCU flow & pool, offer, reservation, issuance, redemption & reconciled simulated accounting \\
Four domains & system, redistribution, user-proxy, accounting & раздельные evaluation outputs \\
Joint robustness & Wilson rule, threshold, not-evaluable handling & scenario-conditional decision rule \\
Decision & pilot evidence и partner ask & следующий проверяемый шаг \\
\bottomrule
\end{longtable}

\textbf{Главный демонстрационный путь:} один реальный Vilnius site/source/timestamp $\to$ qualified или withheld scarcity state $\to$ bounded OSA signal $\to$ funded alternative scenario $\to$ assumed response $\to$ decision boundary. Если data gate закрыт, formal effectiveness остаётся locked, но можно показать отдельный coloured preview с постоянной надписью \emph{Illustrative scenario --- current history is not yet data-qualified}. Цвет кодирует только рассчитанное изменение release probability при declared assumptions, а не доказанную эффективность policy.

\section{Как проводить demo}
\begin{enumerate}
\item \textbf{Evidence:} покажите карту, реальные состояния, tariff structure и readiness gate. Скажите, что большое число строк не заменяет календарную длительность.
\item \textbf{Policy:} покажите origin signal $\rightarrow$ MCU pool $\rightarrow$ funded receiver offer; отделите offer от issuance.
\item \textbf{Response:} покажите price, credit и distance sensitivities как assumptions.
\item \textbf{Evaluation:} покажите четыре domain cards и not-evaluable boundaries.
\item \textbf{Robustness:} покажите illustrative preview и одним предложением отделите его от будущего data-qualified Wilson result.
\item \textbf{Decision:} завершите bounded pilot и поясните, что DAEI --- future pathway, не текущий MCU.
\end{enumerate}
Никогда не тратьте demo на прокрутку кода. Код нужен для ответа о reproducibility; интерфейс --- для рассказа решения.

\section{Трёхминутный pitch: смысловой каркас}
\textbf{0:00--0:25.} «Вильнюс видит состояние зарядок, но не видит очередь и не знает, как водители ответят на новую цену. Поэтому обещать эффект было бы нечестно.»

\textbf{0:25--1:00.} «Мы превратили OCPI и DATEX в единый evidence pipeline: восстановили доступность, интервалы и scarcity, а реальный OCPI тариф используем как baseline.»

\textbf{1:00--1:30.} «OSA-v1 определяет bounded scarcity signal, а prior funded MCU поддерживают qualified alternatives. Offer, issuance и redemption не смешиваются; MCU не является DAEI.»

\textbf{1:30--2:20.} «Поскольку response неизвестен, мы варьируем price, credit, distance и delay assumptions. Fixed и adaptive используют общий path и random numbers; choice проходит safety и reconciled ledger.»

\textbf{2:20--2:45.} «Результат --- four-domain map. Joint success требует system, redistribution, user-proxy и accounting gates одновременно.»

\textbf{2:45--3:00.} «Наш ask: продолжить сбор, выбрать qualified pilot sites и заменить assumptions измеренными Vilnius effects.»

\boxtext{\textbf{Главная фраза:} We do not claim that the policy is already proven; we identify the joint price-and-credit response conditions under which it is robust within observable limits, expose spatial, access-proxy and accounting trade-offs, and turn them into a measurable city pilot.}

\section{Вопросы жюри}
\textbf{Почему не только pricing?} Scarcity signal остаётся управляемым trigger, но ценовой вклад направляется в balanced MCU circulation для qualified alternatives. Это не замена инфраструктуре и не production payment system.

\textbf{Это DAEI?} Нет. DAEI --- регулируемая renewable-transport accounting system и возможный future integration precedent. MCU --- отдельная simulated charging-balancing unit без conversion rate к DAEI.

\textbf{Откуда вы знаете реакцию людей?} Пока не знаем. Поэтому $\theta$ маркирован assumed, выполняется sensitivity analysis, а threshold превращается в измеримую гипотезу пилота.

\textbf{Где очередь?} В данных её нет, и мы её не заявляем. Zero availability --- строго определённый scarcity evidence.

\textbf{Почему вы показываете heatmap до 7 дней?} Formal effectiveness heatmap всё ещё заблокирован. Видимый preview показывает только форму assumed release-response mechanism; он постоянно маркирован как illustrative, не проходит в robustness decision и не выдаётся за эффект на Vilnius data.

\textbf{Вы просто подавляете demand?} Мы отдельно показываем accepted-start displacement, unassigned demand и receiver scarcity. Успех по одной scarcity metric недостаточен.

\textbf{Зачем OCPI и DATEX?} OCPI лучше описывает EVSE/status/tariff, DATEX дополняет инфраструктурный контекст. Мы не смешиваем их без audited mapping.

\textbf{В чём инновация?} Не «магическая формула», а decision-support under sparse behavioural data: реальные тарифы + цифровой двойник + explicit uncertainty + falsifiable pilot threshold.

\textbf{Что сделать городу завтра?} Продолжить collection, выбрать sites после readiness audit, обеспечить Session/CDR, провести bounded A/B или staggered rollout и оценить заранее заданные outcomes.

\textbf{Кто платит?} Рабочая гипотеза --- city--operator pilot partnership. Prior funded contributions и явно раскрытый subsidy могут финансировать MCU scenario; real liability и settlement должны быть согласованы до выдачи value пользователям.

\textbf{Какой один KPI?} Qualified zero-availability snapshot share должен снизиться без ухудшения receiver scarcity; accounting при этом остаётся solvent and reconciled.

\section{Подготовка команды}
\textbf{Сессия 1 (60 мин):} вместе нарисовать pipeline по памяти; объяснить друг другу provenance и claim scopes.\\
\textbf{Сессия 2 (60 мин):} один объясняет F1/OSA, второй --- transition/response/robustness; затем поменяться.\\
\textbf{Сессия 3 (45 мин):} demo без текста; зафиксировать backup screenshots и offline story.\\
\textbf{Сессия 4 (45 мин):} два pitch подряд с таймером и 15 hostile questions.\\
\textbf{Перед сценой:} проверить collector/database copy, UI, resolution, seed, readiness timestamp; не обещать live network; иметь PDF и 3 ключевых screenshots.

\section{Одностраничная шпаргалка}
\begin{itemize}
\item Product: policy stress-testing digital twin for Vilnius EV charging.
\item Evidence: OCPI + DATEX; EVSE is capacity unit; UTC storage/Vilnius calendar.
\item Problem metric: observable operational scarcity, never an invented queue.
\item Baseline: F1 structured observed tariff replay.
\item Policy: OSA-v1 scarcity signal + prior-funded MCU circulation; offer $\neq$ issuance $\neq$ redemption.
\item Separation: $\lambda$ intensity, $\kappa$ coverage, $R_i(t)$ rollout.
\item Unknown behaviour: explicit price, credit, distance and delay assumptions.
\item Fair comparison: same path, same random numbers, same metrics.
\item Decision: Wilson lower bound $\ge 0.8$, minimum 20 replications.
\item Guardrails: displaced starts, equipment loss, spatial spillover.
\item Current limitation: 4.08 days at the 10 October snapshot; 7-day transition gate not yet passed; preview is illustrative only.
\item Ask: collect $\rightarrow$ qualify sites $\rightarrow$ bounded pilot $\rightarrow$ validate response.
\end{itemize}

\section{Финальная проверка понимания}
Если вы оба можете без подсказки ответить на пять вопросов, вы готовы: (1) почему zero availability не queue; (2) чем F1 отличается от OSA; (3) почему $\theta$ assumed; (4) как paired Monte Carlo делает сравнение честнее; (5) чем illustrative preview отличается от formal effectiveness result.

\vfill
\small\textit{Guide описывает фактическое состояние репозитория и локальной базы на 10 октября 2026. Числа readiness должны пересчитываться перед записью и презентацией, поскольку collector продолжает добавлять данные. Нормативные детали определяются versioned Markdown contracts в docs/.}
\end{document}
